\medskip

\textbf{Question 1}

Quelle est la mesure, en degrés, d'un angle droit ?

\medskip

\textbf{Question 2}

Voici une série de quatre notes : 8~;~10~;~11~;~11. 

Quelle est la moyenne de cette série ?

\begin{center}
\begin{tabularx}{\linewidth}{*{4}{X}}
\textbf{A.~~}9,5 &\textbf{B.~~}10 &\textbf{C.~~}10,5 &\textbf{D.~~}11
\end{tabularx}
\end{center}
\medskip

\textbf{Question 3}


Dans un collège de 800 élèves, 25\,\% des élèves portent des lunettes.

Combien d'élèves portent des lunettes ?



\medskip

\textbf{Question 4}

Le graphique ci-dessous donne l'évolution de la température (en degrés Celsius) en fonction de l'horaire (en heures).

\begin{center}
\psset{xunit=0.5cm,yunit=0.125cm}
\begin{pspicture}(-2,-5)(18,35)
\multido{\n=0+2}{10}{\psline[linewidth=0.2pt](\n,0)(\n,35)}
\multido{\n=0+5}{8}{\psline[linewidth=0.2pt](0,\n)(18,\n)}
\psaxes[linewidth=1.25pt,Dx=2,Dy=5,labelFontSize=\scriptstyle](0,0)(0,0)(18,35)
\psdots[dotscale=1.5](8,15)(12,27)(16,30)
\psline(8,15)(12,27)(16,30)
\uput[d](9,-3){Horaire (en heures)}
\rput{90}(-2,17.5){Température (en $^\circ$C)}
\end{pspicture}
\end{center}

Entre 8~h et 16~h, de combien de degrés la température a-t-elle augmenté ?

\medskip

\textbf{Question 5}

Une voiture roule à 90~km/h. Combien de temps met-elle pour parcourir 45~km ?~
\begin{center}
\begin{tabularx}{\linewidth}{*{4}{X}}
\textbf{A.~~}15 min &\textbf{B.~~}30 min &\textbf{C.~~}45~min &\textbf{D.~~}1~h
\end{tabularx}
\end{center}

\medskip

\begin{minipage}{0.65\linewidth}
\textbf{Question 6}

Donner le périmètre du losange ABCD représenté ci- contre.
\end{minipage}\hfill
\begin{minipage}{0.32\linewidth}
\psset{unit=1cm}
\begin{pspicture}(-2.5,-2.2)(2.5,2.2)
\uput[l](-2.3,0){A} \uput[u](0,2){D} \uput[r](2.3,0){C} \uput[d](0,-2){B} \uput[ul](-1.15,1){3 cm} 
\pspolygon(-2.3,0)(0,2)(2.3,0)(0,-2)%ABCD
\end{pspicture}
\end{minipage}

\medskip

\textbf{Question 7} (1 point)

Pour résoudre l'équation $4x - 3 = 0$, on effectue le calcul :

\begin{center}
\begin{tabularx}{\linewidth}{*{4}{X}}
\textbf{A.~~}$x = \dfrac{20}{4} + 3$ &\textbf{B.~~}$x = (20 - 4) + 3$ &\textbf{C.~~}$x = 20 \times 4 + 3$ &\textbf{D.~~}$x = \dfrac{20 + 3}{4}$
\end{tabularx}
\end{center}

\medskip

\textbf{Question 8} (1 point)

\begin{minipage}{0.55\linewidth}
Sur la figure ci-contre, les droites (DE) et (AC) sont parallèles.

Écrire une égalité de rapports permettant de déterminer la longueur AB.

\end{minipage}\hfill
\begin{minipage}{0.4\linewidth}
\psset{unit=1cm}
\begin{pspicture}(5.4,4)
\pspolygon(0.3,0.1)(4.7,0.6)(0.4,3.4)%ACB
\psline(0.35,1.2)(3.37,1.5)%DE
\uput[dl](0.3,0.1){A}  \uput[r](4.7,0.6){C} \uput[ul](0.4,3.4){B} 
\uput[l](0.35,1.2){D}\uput[ur](3.37,1.5){E} 
\uput[d](2.5,0.35){6}\uput[d](1.86,1.35){4}\uput[ur](-0.48,1.85){4,5}
\end{pspicture}
\end{minipage}
\medskip

\textbf{Question 9}

\begin{minipage}{0.47\linewidth}
On considère l'algorithme suivant :

Quel résultat obtient-on si on choisit 1 comme nombre de départ ?
\end{minipage}\hfill
\begin{minipage}{0.5\linewidth}
\begin{scratch}[scale=0.8]
\blockinit{quand \greenflag est cliqué}
\blocksensing{demander \ovalnum{Choisir un nombre} et attendre}
\blockvariable{mettre \selectmenu{\ovalvariable{variable}} à \ovalsensing{réponse}}
\blockvariable{mettre \selectmenu{résultat} à \ovaloperator{\ovalnum{8} *  \ovalvariable{variable}	}}
\blockvariable{mettre \selectmenu{résultat} à \ovaloperator{\ovalvariable{résultat} + \ovalnum{10}}}
\blockvariable{mettre \selectmenu{résultat} à \ovaloperator{\ovalvariable{résultat} / \ovalnum{2}}}
\blocklook{dire \ovaloperator{regrouper \ovalnum{J'obtiens comme résultat} et \ovalvariable{résultat}}}
\end{scratch}
\end{minipage}

\medskip
